\section{New Capabilities}
\label{sec:capabilities}

\subsection{Building Blocks}
\label{subsec:blocks}

Long-running multi-agent systems typically share many basic structural attributes. The particular mechanisms by which each of these is implemented is what enables a system to excel at its target use case.
\begin{enumerate}
	\item Context window management: dealing with LLM sessions, optimizing for prefix hit rates, operating around context limits
	\item Subagent session management and control: structuring the tasks to get high-quality results
	\item Memory: maintaining state across multiple LLM sessions and over the duration of a project that would otherwise take a human team months
	\item Knowledge: many systems incorporate specialized knowledge or skills through various mechanisms
\end{enumerate}

In Design Conductor 2.0, all of these attributes were rebuilt to support longer-horizon and more complex semiconductor design tasks. 

\subsection{Novel Capabilities}

Based on the building blocks described in \ref{subsec:blocks}, and paired with models released in April 2026, Design Conductor 2.0 acquired several new capabilities that made it possible to build the designs described in section \ref{sec:designs}.

\subsubsection{Concept to Layout}

In contrast to the original Design Conductor 1.0, the new version is able to develop its own specifications. It is able to take a concept -- like a TurboQuant accelerator -- and autonomously map it to a performant design. This means that it can actually be a true partner to master designers, and that the right interface is no longer spec to RTL or spec to layout, but rather concept onwards. Design Conductor 2.0 can relieve the team from much of the work of building exhaustive and authoritative specifications.

\subsubsection{Architecture understanding \& first-principles reasoning}

One of the primary areas where Design Conductor 1.0 struggled was in making reasoned architecture and system design tradeoffs. This constrained its applicability in more complex designs where engineers need to make these tradeoffs on a regular basis. The new harness effectively solves this problem.

For example, in the case of VerTQ, it developed a custom exponentiation unit that uses a (to our knowledge) novel factoring of a fifth-degree polynomial to minimize delay at the precision the task required. It considered how to structure and map the VerTQ design to the VU29P multi-die FPGA, minimizing communication overhead. This ability to reason from first principles is what makes it possible for Design Conductor 2.0 to tackle completely new problems with innovative solutions.

\subsubsection{Closing the loop}

Thanks to the other capabilities noted above, Design Conductor 2.0 has the ability to ``close the loop'', taking high-level feedback from timing analysis and placement tools, reasoning about it, and making complex RTL changes to adapt to this. In the case of VerTQ, we tested this by modifying the timing target once an initial version was built. When this happens, Conductor refactors large parts of the design to achieve the PPA goals. This blunts the impact of one of the most painful experiences for an all-human design team -- finding out that a design isn't going to meet timing and needs to be restructured.

\subsubsection{Mathematical Reasoning Abilities}

Conductor leverages the improved mathematical reasoning abilities of the new models. In particular, Conductor 2.0 is able to accurately analyze design characteristics, as well as to optimize the implementation of numerics for specific applications.
